\begin{document}
\maketitle

\clearpage

\section*{Resumo}\label{resumo}
\addcontentsline{toc}{section}{Resumo}

Este trabalho investiga o impacto de diferentes estratégias de paralelização em Unidades de Processamento Gráfico (GPUs) no desempenho de um algoritmo memético (Algoritmo Genético híbrido com busca local 2-opt) para o Problema do Caixeiro Viajante (TSP).
Foram implementadas e comparadas quatro variantes de um GA memético: uma versão sequencial em CPU, uma versão híbrida ingênua (transferência indivíduo a indivíduo), uma versão híbrida otimizada (processamento em lote) e uma versão totalmente residente em GPU.
Os experimentos, conduzidos em um conjunto de 38 instâncias da TSPLIB com até 1002 cidades, demonstram que a estratégia de paralelização influencia drasticamente o trade-off entre tempo de execução e qualidade da solução. A variante híbrida otimizada obteve os menores tempos de execução (ganho de performance mediano por instância de 5,69\times sobre a híbrida ingênua), enquanto a variante totalmente em GPU alcançou a melhor qualidade de solução (gap médio de 0,88\%).
Os resultados evidenciam que a minimização da transferência de dados entre CPU e GPU é crítica para o desempenho, e que arquiteturas totalmente residentes favorecem o refinamento das soluções em problemas de otimização combinatória.

\textbf{Ressalva:} os valores numéricos citados neste resumo são provisórios e não devem ser lidos como resultados finais; ver a \hyperref[obs-ajustes]{Seção 1.3}.

\textbf{Palavras-chave:} Problema do Caixeiro Viajante, Algoritmos Genéticos, Computação em GPU, CUDA, Otimização Combinatória.

\clearpage

\section*{Abstract}\label{abstract}
\addcontentsline{toc}{section}{Abstract}

This work investigates the impact of different parallelization strategies on Graphics Processing Units (GPUs) on the performance of a memetic algorithm (Genetic Algorithm hybridized with 2-opt local search) for the Traveling Salesman Problem (TSP).
Four analogous variants of a memetic GA were implemented and compared: a sequential CPU version, a naive hybrid version (individual-by-individual transfer), an optimized hybrid version (batch processing), and a fully GPU-resident version.
Experiments conducted on a set of 38 TSPLIB instances with up to 1002 cities demonstrate that the parallelization strategy drastically influences the trade-off between execution time and solution quality.
The optimized hybrid variant achieved the lowest execution times (median per-instance Performance gains of \(5.69\times\) over the naive hybrid), while the fully GPU-resident variant achieved the best solution quality (average gap of \(0.88\%\)).
The results highlight that minimizing data transfer between CPU and GPU is critical for performance, and that fully resident architectures favor solution refinement in combinatorial optimization problems.

\textbf{Note:} the numerical values cited in this abstract are provisional and should not be read as final results; see \hyperref[obs-ajustes]{Section 1.3}.

\textbf{Keywords:} Traveling Salesman Problem, Genetic Algorithms, GPU
Computing, CUDA, Combinatorial Optimization.

\clearpage

\clearpage
\csname tableofcontents\endcsname
\clearpage

\section*{Glossário de termos}\label{glossario}
\addcontentsline{toc}{section}{Glossário de termos}

\begin{longtable}[]{@{}
    >{\raggedright\arraybackslash}p{(\linewidth - 4\tabcolsep) * \real{0.3333}}
    >{\raggedright\arraybackslash}p{(\linewidth - 4\tabcolsep) * \real{0.3333}}
    >{\raggedright\arraybackslash}p{(\linewidth - 4\tabcolsep) * \real{0.3333}}@{}}
  \toprule\noalign{}
  \begin{minipage}[b]{\linewidth}\raggedright
    Termo em inglês
  \end{minipage} & \begin{minipage}[b]{\linewidth}\raggedright
                     Forma adotada no texto
                   \end{minipage}                   & \begin{minipage}[b]{\linewidth}\raggedright
                                                        Observação
                                                      \end{minipage}                                                                 \\
  \midrule\noalign{}
  \endhead
  \bottomrule\noalign{}
  \endlastfoot
  \emph{batch}                                & lote                                                          & Usado para processamento de vários indivíduos em
  uma única chamada.                                                                                                                                             \\
  \emph{benchmark}                            & benchmark / avaliação experimental                            & Mantém-se
  ``benchmark'' quando o foco é o protocolo comparativo.                                                                                                         \\
  \emph{crossover}                            & cruzamento                                                    & Preferência por ``cruzamento'', mantendo
  a sigla OX quando necessário.                                                                                                                                  \\
  \emph{fitness}                              & aptidão                                                       & Quando útil, escreve-se ``aptidão
  (\emph{fitness})'' na primeira ocorrência.                                                                                                                     \\
  \emph{full GPU} / \emph{GPU-resident}       & totalmente em GPU / residente na
  GPU                                         & Expressões usadas para a variante com dados mantidos em VRAM.                                                    \\
  \emph{gap}                                  & gap / desvio percentual em relação ao ótimo                   & O termo
  ``gap'' é preservado por ser consagrado na literatura.                                                                                                         \\
  \emph{host}                                 & hospedeiro (CPU)                                              & Contraposto a \emph{device} (GPU).               \\
  \emph{kernel launch}                        & lançamento de \emph{kernel}                                   & O termo
  \emph{kernel} é mantido por ser técnico e específico de CUDA.                                                                                                  \\
  \emph{offspring}                            & descendentes                                                  & Refere-se aos indivíduos gerados após
  cruzamento e mutação.                                                                                                                                          \\
  \emph{overhead}                             & sobrecarga                                                    & Usado para custos adicionais de
  comunicação, sincronização ou inicialização.                                                                                                                   \\
  \emph{speedup}                              & ganho de performance                                          & Preferência por esta expressão
  no texto analítico.                                                                                                                                            \\
  \emph{swap mutation}                        & mutação por troca                                             & Operador que troca a posição
  de duas cidades na permutação.                                                                                                                                 \\
  \emph{throughput}                           & vazão                                                         & Preferência por ``vazão'', mantendo o inglês
  quando a literatura exigir.                                                                                                                                    \\
  \emph{warp}                                 & \emph{warp} (grupo de 32 \emph{threads})                      & O termo inglês
  é mantido, com explicação curta na primeira ocorrência.                                                                                                        \\
\end{longtable}

\clearpage

\section*{Lista de Símbolos}\label{lista-de-simbolos}
\addcontentsline{toc}{section}{Lista de Símbolos}

\begin{longtable}[]{@{}
    >{\raggedright\arraybackslash}p{(\linewidth - 4\tabcolsep) * \real{0.3333}}
    >{\raggedright\arraybackslash}p{(\linewidth - 4\tabcolsep) * \real{0.3333}}
    >{\raggedright\arraybackslash}p{(\linewidth - 4\tabcolsep) * \real{0.3333}}@{}}
  \toprule\noalign{}
  \begin{minipage}[b]{\linewidth}\raggedright
    Símbolo
  \end{minipage} & \begin{minipage}[b]{\linewidth}\raggedright
                     Descrição
                   \end{minipage}             & \begin{minipage}[b]{\linewidth}\raggedright
                                                  Ver seção
                                                \end{minipage}                \\
  \midrule\noalign{}
  \endhead
  \bottomrule\noalign{}
  \endlastfoot
  \(n_{coords}\)                              & Número de cidades (coordenadas) da instância            &
  \hyperref[sec-33-selecao-instancias]{3.3}                                                               \\
  \(n_{pop}\)                                 & Tamanho da população do Algoritmo Genético              &
  \hyperref[sec-321-ga-base]{3.2.1}                                                                       \\
  \(n_{reps}\)                                & Número de repetições independentes por experimento      &
  \hyperref[sec-33-selecao-instancias]{3.3}                                                               \\
  \(p_{\text{mut}}\)                          & Taxa de mutação                                         &
  \hyperref[sec-23-algoritmos-geneticos]{2.3}                                                             \\
  \(p_c\)                                     & Taxa de cruzamento                                      &
  \hyperref[sec-23-algoritmos-geneticos]{2.3}                                                             \\
  \(k_{\text{trnmt}}\)                        & Tamanho do torneio na seleção                           &
  \hyperref[sec-251-ga-2opt]{2.5.1}                                                                       \\
  \(n_{2\text{-opt}}\)                        & Número máximo de iterações da busca local 2-opt         &
  \hyperref[sec-251-ga-2opt]{2.5.1}                                                                       \\
  \(B_{\mathrm{H2D}}\)                        & Volume de dados transferidos Host-to-Device             &
  \hyperref[cap3-materiais-metodos]{Cap. 3}                                                               \\
  \(B_{\mathrm{D2H}}\)                        & Volume de dados transferidos Device-to-Host             &
  \hyperref[cap3-materiais-metodos]{Cap. 3}                                                               \\
  \(T\)                                       & Tempo de execução                                       &
  \hyperref[sec-41-resultados-experimentais]{4.1}                                                         \\
  \(\Delta C\)                                & Variação de custo em um movimento local                 &
  \hyperref[sec-221-2opt]{2.2.1}                                                                          \\
  \(G = (V, E)\)                              & Grafo composto por vértices \(V\) e arestas \(E\)       &
  \hyperref[sec-21-otimizacao-tsp]{2.1}                                                                   \\
  \(c_{ij}\)                                  & Custo (distância) entre as cidades \(i\) e \(j\)        &
  \hyperref[sec-21-otimizacao-tsp]{2.1}                                                                   \\
  \(f(x)\)                                    & Função de aptidão (\emph{fitness}) de uma solução \(x\) &
  \hyperref[sec-23-algoritmos-geneticos]{2.3}                                                             \\
  \(P(t)\)                                    & População de soluções na geração \(t\)                  &
  \hyperref[sec-23-algoritmos-geneticos]{2.3}                                                             \\
  \(\kappa\)                                  & Limite de iterações para busca local                    &
  \hyperref[sec-221-2opt]{2.2.1} e \hyperref[sec-251-ga-2opt]{2.5.1}                                      \\
\end{longtable}

\clearpage

\listoffigures

\clearpage

\listoftables

